\section{Pandas}\label{sec:07-pandas}

\subsection{DataFrame erzeugen und darauf zugreifen}\label{sec:07-dataframe-erzeugen}
\begin{contentbox}
\CodeLine{import pandas as pd}\\
\CodeLine{df = pd.DataFrame()}[leerer DataFrame]\\
\CodeLine{df = pd.read_csv("Name.csv")}[DataFrame aus CSV laden]\\
\ZSFgapS
\OverlineComment{Zugriff auf Daten}
\CodeLine{data["Month"]}[einzelne Spalte]\\
\CodeLine{data[["Name", "City"]]}[mehrere Spalten]\\
\CodeLine{data["Name"][2]}[gibt `Carlos' zurück]\\
\ZSFgapS
\CodeLine{data.iloc[1]}[Zeile mit Index 1]\\
\CodeLine{data[1:3]}[Zeilen mit Index 1 und 2]\\
\CodeLine{data.loc[1:3, "M":"C"]}[Zeilen 1--3, Spalten M bis C]\\
\CodeLine{data.iloc[1:3, 0:2]}[wie loc, aber indexbasiert]
\end{contentbox}


\subsection{Spalten umbenennen}\label{sec:07-spalten-umbenennen}
\begin{contentbox}
\CodeLine{data.set_index("Row")}[setzt `Row' als Index]\\
\CodeLine{data.rename(columns={"City": "Location"})}[Spalte umbenennen]\\
\CodeLine{data.columns = ["Monat", "Name", "Stadt", ...]}[jede Spalte umbenennen] % chktex 11 chktex 26
\end{contentbox}

\subsection{Filtern von Zeilen}\label{sec:07-filtern-von-zeilen}
\begin{contentbox}
\CodeLine{data[data["Age"] > 30]}[alle Personen über 30]\\
\CodeLine{data["Name"][data["Age"] > 30]}[nur Namen dieser Personen]
\end{contentbox}


\subsection{Zeilen und Spalten löschen}\label{sec:07-zeilen-und-spalten-loeschen}
\begin{contentbox}
\CodeLine{data.drop(columns=["Age"])}[löscht Spalte `Age']\\
\CodeLine{data.drop(data.index[0])}[löscht erste Zeile]\\
\ZSFgapS
\OverlineComment{löscht Zeilen mit NaN}
\CodeLine{data.dropna(axis=0, how="any")}[wenn mind. 1 NaN in der Zeile]\\
\CodeLine{data.dropna(axis=0, how="all")}[wenn alle NaN in der Zeile]\\
\ZSFgapS
\CodeLine{data.dropna(axis=1, how="any")}[löscht Spalten mit mind. 1 NaN]
\end{contentbox}


\subsection{Hilfsfunktionen für DataFrames}\label{sec:07-hilfsfunktionen-fuer-dataframes}
\begin{contentbox}
\CodeLine{data["Age"].sum()}[summiert alle Werte der Spalte `Age']\\
\CodeLine{data["Age"].max()}[maximaler Wert]\\
\CodeLine{data.fillna(...)}[ersetzt NaN durch gewünschten Wert] % chktex 11
\end{contentbox}






%=====================================================================================================================
%%%%%%%%%%%%%%%%%%%%%%%%%%%%%%%%%%%%%%%
